% Part: first-order-logic
% Chapter: syntax-and-semantics
% Section: assignments

\documentclass[../../../include/open-logic-section]{subfiles}

\begin{document}

\olfileid[tr]{fol}{syn}{ass}

\olsection{Değişken Atamaları}

\begin{explain}
Bir !!{variable} ataması~$s$, \emph{her} değişkene bir değer verir; bu
değişkenlerden sonsuz sayıda vardır. Elbette buna gerek yoktur.
!!{variable} atamalarının bütün !!{variable}s ifadelerine değer
atamasını yalnızca işleri çok kolaylaştırdığı için isteriz. Bir
terim~$t$'nin değeri ve !!a{formula}~$!A$ ifadesinin !!a{structure}
içinde~$s$'ye göre sağlanıp sağlanmaması, yalnızca~$s$'nin~$t$ içindeki
!!{variable}s ifadelerine ve~$!A$'nın serbest !!{variable}s ifadelerine
yaptığı atamalara bağlıdır. Sonraki iki önerme tam olarak bunu söyler.
``Bağlıdır'' düşüncesini kesinleştirmek için,~$t$ içindeki bütün
değişkenler üzerinde uyuşan herhangi iki değişken atamasının aynı değeri
verdiğini ve~$!A$'nın bütün serbest değişkenleri üzerinde uyuşan iki
değişken ataması varsa~$!A$'nın birine göre sağlanmasının ancak ve ancak
ötekine göre sağlanmasıyla eşdeğer olduğunu gösteririz.
\end{explain}

\begin{prop}\ollabel{prop:valindep}
Bir terim~$t$ içindeki !!{variable}s, $x_1$, \dots,~$x_n$ arasındaysa
ve $i = 1$, \dots,~$n$ için $s_1(x_i) = s_2(x_i)$ ise
$\Value{t}{M}[s_1] = \Value{t}{M}[s_2]$.
\end{prop}

\begin{proof}
~$t$'nin karmaşıklığı üzerinde tümevarımla. Temel durumda~$t$,
!!a{constant} ya da $x_1$, \dots,~$x_n$ değişkenlerinden biri olabilir.
$t = c$ ise $\Value{t}{M}[s_1] = \Assign{c}{M} =
\Value{t}{M}[s_2]$. $t = x_i$ ise önerme varsayımı gereği
$s_1(x_i) = s_2(x_i)$; dolayısıyla $\Value{t}{M}[s_1] = s_1(x_i) =
s_2(x_i) = \Value{t}{M}[s_2]$.

Tümevarım adımı için $t = \Atom{f}{t_1, \dots, t_k}$ olduğunu ve savın
$t_1$, \dots, $t_k$ için geçerli olduğunu varsayalım. O zaman
\begin{align*}
  \Value{t}{M}[s_1] & = \Value{\Atom{f}{t_1, \dots, t_k}}{M}[s_1] \\
  & = \Assign{f}{M}(\Value{t_1}{M}[s_1], \dots, \Value{t_k}{M}[s_1]).
\intertext{$j = 1$, \dots,~$k$ için~$t_j$ içindeki !!{variable}s,
    $x_1$, \dots,~$x_n$ arasındadır. Tümevarım varsayımı gereği
    $\Value{t_j}{M}[s_1] = \Value{t_j}{M}[s_2]$. Dolayısıyla}
\Value{t}{M}[s_1] & = \Value{\Atom{f}{t_1, \dots, t_k}}{M}[s_1] \\
  & = \Assign{f}{M}(\Value{t_1}{M}[s_1], \dots, \Value{t_k}{M}[s_1]) \\
  & = \Assign{f}{M}(\Value{t_1}{M}[s_2], \dots, \Value{t_k}{M}[s_2]) \\
  & = \Value{\Atom{f}{t_1, \dots, t_k}}{M}[s_2] = \Value{t}{M}[s_2].
\end{align*}
\end{proof}

\begin{prop}\ollabel{prop:satindep}
$!A$ içindeki serbest !!{variable}s, $x_1$, \dots,~$x_n$ arasındaysa
ve $i = 1$, \dots,~$n$ için $s_1(x_i) = s_2(x_i)$ ise
$\Sat{M}{!A}[s_1]$ ancak ve ancak $\Sat{M}{!A}[s_2]$.
\end{prop}

\begin{proof}
$!A$'nın karmaşıklığı üzerinde tümevarım kullanırız. $!A$'nın atomik
olduğu temel durumda~$!A$ şunlardan biri olabilir:
\iftag{prvTrue}{$\ltrue$,}{}
\iftag{prvFalse}{$\lfalse$,}{}
$k$-li bir yüklem~$R$ ve $t_1$, \dots,~$t_k$ terimleri için
$\Atom{R}{t_1, \dots, t_k}$ ya da $t_1$ ve~$t_2$ terimleri için
$\eq[t_1][t_2]$. Son iki durumda !!{biconditional} ifadesinin yalnızca
ileri yönünü gösteririz; çünkü ters yönün ispatı simetriktir.

\begin{enumerate}
\tagitem{prvTrue}{%
  \indcase{!A}{\ltrue}{hem $\Sat{M}{\indfrm}[s_1]$ hem de
    $\Sat{M}{\indfrm}[s_2]$.}}{}

\tagitem{prvFalse}{%
  \indcase{!A}{\lfalse}{hem $\Sat/{M}{!A}[s_1]$ hem de
    $\Sat/{M}{!A}[s_2]$.}}{}

\item
  \indcase{!A}{\Atom{R}{t_1, \ldots, t_k}}{
    $\Sat{M}{\indfrm}[s_1]$ olsun. O zaman
    \[
    \langle \Value{t_1}{M}[s_1], \ldots, \Value{t_k}{M}[s_1] \rangle
    \in \Assign{R}{M}.
    \]
    $i = 1$, \dots,~$k$ için \olref{prop:valindep} gereği
    $\Value{t_i}{M}[s_1] = \Value{t_i}{M}[s_2]$. Dolayısıyla ayrıca
    $\langle \Value{t_1}{M}[s_2], \ldots, \Value{t_k}{M}[s_2] \rangle
    \in \Assign{R}{M}$ ve buradan $\Sat{M}{\indfrm}[s_2]$.}
\item
  \indcase{!A}{\eq[t_1][t_2]}{$\Sat{M}{\indfrm}[s_1]$ varsayılsın.
    O zaman $\Value{t_1}{M}[s_1] = \Value{t_2}{M}[s_1]$. Dolayısıyla
    \begin{align*}
      \Value{t_1}{M}[s_2] & = \Value{t_1}{M}[s_1]
      & \text{(\olref{prop:valindep} gereği)} \\
      & = \Value{t_2}{M}[s_1]
      & \text{($\Sat{M}{\eq[t_1][t_2]}[s_1]$ olduğundan)}\\
      &= \Value{t_2}{M}[s_2]
      & \text{(\olref{prop:valindep} gereği),}
    \end{align*}
    öyleyse $\Sat{M}{\eq[t_1][t_2]}[s_2]$.}
\end{enumerate}

Şimdi~$!A$'dan daha az karmaşık bütün !!{formula}s~$!B$ için
$\Sat{M}{!B}[s_1]$ ancak ve ancak $\Sat{M}{!B}[s_2]$ olduğunu
varsayalım. Tümevarım adımı,~$!A$'nın ana işlecinin belirlediği
durumlara göre ilerler. Her durumda !!{biconditional} ifadesinin yalnızca
ileri yönünü gösteririz; ters yönün ispatı simetriktir. Niceleyici
durumları dışındaki bütün durumlarda tümevarım varsayımını~$!A$'nın
alt-!!{formula}s ifadeleri~$!B$'ye uygularız. $!B$'nin serbest
değişkenleri~$!A$'nın serbest değişkenleri arasındadır. Dolayısıyla
$s_1$ ile $s_2$,~$!A$'nın serbest değişkenleri üzerinde uyuşuyorsa
~$!B$'ninkiler üzerinde de uyuşur ve tümevarım varsayımı~$!B$'ye
uygulanır.

\begin{enumerate}
\tagitem{defNot}{}{%
  \iftag{probNot}{%
    \indcase!{!A}{\lnot !B}{}}{%
    \indcase{!A}{\lnot !B}{$\Sat{M}{\indfrm}[s_1]$ ise
      $\Sat/{M}{!B}[s_1]$; tümevarım varsayımı gereği
      $\Sat/{M}{!B}[s_2]$ ve dolayısıyla $\Sat{M}{\indfrm}[s_2]$.}}}

\tagitem{defAnd}{}{%
  \iftag{probAnd}{%
    \indcase!{!A}{!B \land !C}{}}{%
    \indcase{!A}{!B \land !C}{$\Sat{M}{\indfrm}[s_1]$ ise
      $\Sat{M}{!B}[s_1]$ ve $\Sat{M}{!C}[s_1]$; tümevarım varsayımı
      gereği $\Sat{M}{!B}[s_2]$ ve $\Sat{M}{!C}[s_2]$. Dolayısıyla
      $\Sat{M}{\indfrm}[s_2]$.}}}

\tagitem{defOr}{}{%
  \iftag{probOr}{%
    \indcase!{!A}{!B \lor !C}{}}{%
    \indcase{!A}{!B \lor !C}{$\Sat{M}{\indfrm}[s_1]$ ise
      $\Sat{M}{!B}[s_1]$ ya da $\Sat{M}{!C}[s_1]$. Tümevarım
      varsayımı gereği $\Sat{M}{!B}[s_2]$ ya da $\Sat{M}{!C}[s_2]$;
      dolayısıyla $\Sat{M}{\indfrm}[s_2]$.}}}

\tagitem{defIf}{}{%
  \iftag{probIf}{%
    \indcase!{!A}{!B \lif !C}{}}{%
    \indcase{!A}{!B \lif !C}{$\Sat{M}{\indfrm}[s_1]$ ise
    $\Sat/{M}{!B}[s_1]$ ya da $\Sat{M}{!C}[s_1]$. Tümevarım varsayımı
    gereği $\Sat/{M}{!B}[s_2]$ ya da $\Sat{M}{!C}[s_2]$; dolayısıyla
    $\Sat{M}{\indfrm}[s_2]$.}}}

\tagitem{defIff}{}{%
  \iftag{probIff}{%
    \indcase!{!A}{!B \liff !C}{}}{%
    \indcase{!A}{!B \liff !C}{$\Sat{M}{\indfrm}[s_1]$ ise ya
      $\Sat{M}{!B}[s_1]$ ve $\Sat{M}{!C}[s_1]$, ya da
      $\Sat/{M}{!B}[s_1]$ ve $\Sat/{M}{!C}[s_1]$. Tümevarım varsayımı
      gereği ya $\Sat{M}{!B}[s_2]$ ve $\Sat{M}{!C}[s_2]$, ya da
      $\Sat/{M}{!B}[s_2]$ ve $\Sat/{M}{!C}[s_2]$. Her iki durumda da
      $\Sat{M}{\indfrm}[s_2]$.}}}

\tagitem{defEx}{}{%
  \iftag{probEx}{%
    \indcase!{!A}{\lexists[x][!B]}{}}{%
    \indcase{!A}{\lexists[x][!B]}{$\Sat{M}{\indfrm}[s_1]$ ise
      $\Sat{M}{!B}[\Subst{s_1}{m}{x}]$ olacak bir
      $m \in \Domain{M}$ vardır. $s_1' = \Subst{s_1}{m}{x}$ ve
      $s_2' = \Subst{s_2}{m}{x}$ olsun. $!B$'nin serbest değişkenleri
      $x_1$, \dots, $x_n$ ve $x$ arasındadır. $s_1'$ ile $s_2'$ sırasıyla
      $s_1$ ve~$s_2$'nin $x$-varyantları olduğundan ve varsayım gereği
      $s_1(x_i) = s_2(x_i)$ olduğundan $s_1'(x_i) = s_2'(x_i)$.
      $s_1'$ ve~$s_2'$ tanımlarımız gereği $s_1'(x) = s_2'(x) = m$.
      O halde tümevarım varsayımı~$!B$ ile $s_1'$, $s_2'$ için uygulanır
      ve $\Sat{M}{!B}[s_2']$ elde edilir. $s_2' =
      \Subst{s_2}{m}{x}$ olduğuna göre $\Sat{M}{!B}[\Subst{s_2}{m}{x}]$
      olacak bir $m \in \Domain{M}$ vardır; dolayısıyla
      $\Sat{M}{\indfrm}[s_2]$.}}}

\tagitem{defAll}{}{%
  \iftag{probAll}{%
    \indcase!{!A}{\lforall[x][!B]}{}}{%
    \indcase{!A}{\lforall[x][!B]}{$\Sat{M}{\indfrm}[s_1]$ ise her
      $m \in \Domain{M}$ için $\Sat{M}{!B}[\Subst{s_1}{m}{x}]$.
      Her $m \in \Domain{M}$ için ayrıca
      $\Sat{M}{!B}[\Subst{s_2}{m}{x}]$ olduğunu göstermek istiyoruz.
      Öyleyse gelişigüzel bir $m \in \Domain{M}$ alıp
      $s_1' = \Subst{s_1}{m}{x}$ ve $s_2' = \Subst{s_2}{m}{x}$
      atamalarını ele alalım. $\Sat{M}{!B}[s_1']$. $!B$'nin serbest
      değişkenleri $x_1$, \dots, $x_n$ ve $x$ arasındadır. $s_1'$ ile
      $s_2'$ sırasıyla $s_1$ ve~$s_2$'nin $x$-varyantları olduğundan ve
      varsayım gereği $s_1(x_i) = s_2(x_i)$ olduğundan
      $s_1'(x_i) = s_2'(x_i)$. $s_1'$ ve~$s_2'$ tanımlarımız gereği
      $s_1'(x) = s_2'(x) = m$. O halde tümevarım varsayımı~$!B$ ile
      $s_1'$, $s_2'$ için uygulanır ve $\Sat{M}{!B}[s_2']$ elde edilir.
      Bu, her $m \in \Domain{M}$ için geçerlidir; yani bütün
      $m \in \Domain{M}$ için $\Sat{M}{!B}[\Subst{s_2}{m}{x}]$ ve
      dolayısıyla $\Sat{M}{\indfrm}[s_2]$.}}}
\end{enumerate}
Tümevarımla,~$!A$ içindeki serbest !!{variable}s, $x_1$, \dots, $x_n$
arasındayken ve $i = 1$, \dots,~$n$ için $s_1(x_i)=s_2(x_i)$ iken
$\Sat{M}{!A}[s_1]$ ancak ve ancak $\Sat{M}{!A}[s_2]$ sonucunu elde
ederiz.
\end{proof}

\begin{probtag}{probNot,probOr,probAnd,probIf,probIff,probEx,probAll}
\olref[fol][syn][ass]{prop:satindep} önermesinin ispatını tamamlayın.
\end{probtag}

\begin{explain}
!!^{sentence}s serbest değişken içermez; dolayısıyla herhangi iki
değişken ataması, herhangi bir tümcenin (sıfır tane olan) bütün serbest
değişkenlerine aynı şeyleri atar. Az önce ispatlanan önerme, bu durumda
!!a{sentence} ifadesinin bir yapı içinde değişken atamasına göre
sağlanıp sağlanmamasının atamadan bütünüyle bağımsız olduğu anlamına
gelir. Bu olguyu kaydedelim. Olgu, aşağıda gelecek olan, !!a{sentence}
ifadesinin !!a{structure} içinde sağlanması tanımını (değişken
atamasından söz etmeden) haklı çıkarır.
\end{explain}

\begin{cor}
\ollabel{cor:sat-sentence}
$!A$ !!a{sentence} ve $s$ bir değişken atamasıysa,
$\Sat{M}{!A}[s]$ ancak ve ancak her değişken ataması~$s'$ için
$\Sat{M}{!A}[s']$.
\end{cor}

\begin{proof}
  Gelişigüzel bir değişken ataması~$s'$ alalım. $!A$ !!a{sentence}
  olduğundan serbest değişken içermez; dolayısıyla her değişken
  ataması~$s'$,~$!A$'nın bütün serbest değişkenlerine, önemsiz biçimde,
  ~ $s$ ile aynı şeyleri atar. Böylece \olref{prop:satindep} koşulu
  sağlanır ve $\Sat{M}{!A}[s]$ ancak ve ancak $\Sat{M}{!A}[s']$.
\end{proof}

\begin{defn}
\ollabel{defn:satisfaction}
$!A$ !!a{sentence} ise !!a{structure}~$\Struct M$ yapısının~$!A$'yı
\emph{sağladığını}, yani $\Sat{M}{!A}$ olduğunu, ancak ve ancak bütün
değişken atamaları~$s$ için $\Sat{M}{!A}[s]$ olduğunda söyleriz.
\end{defn}

$\Sat{M}{!A}$ ise kısaca \emph{$!A$,~$\Struct{M}$ içinde doğrudur} da
deriz. Sağlama kavramı, tek tek !!{sentence}s ifadelerinden
!!{sentence}s kümelerine doğal olarak genişler.

\begin{defn}
\ollabel{defn:sat}
$\Gamma$ bir !!{sentence}s kümesiyse, !!a{structure}~$\Struct M$
yapısının~$\Gamma$'yı \emph{sağladığını}, yani $\Sat{M}{\Gamma}$
olduğunu, ancak ve ancak her $!A \in \Gamma$ için $\Sat{M}{!A}$
olduğunda söyleriz.
\end{defn}

\begin{prop}\ollabel{prop:sentence-sat-true}
  $\Struct{M}$ !!a{structure}, $!A$ !!a{sentence} ve $s$ bir değişken
  ataması olsun. $\Sat{M}{!A}$ ancak ve ancak $\Sat{M}{!A}[s]$.
\end{prop}

\begin{proof}
Alıştırma.
\end{proof}

\begin{prob}
\olref[fol][syn][ass]{prop:sentence-sat-true} önermesini ispatlayın.
\end{prob}

\begin{prop}\ollabel{prop:sat-quant}
  $!A(x)$'in yalnızca~$x$'i serbest içerdiğini ve $\Struct M$'nin
  !!a{structure} olduğunu varsayın. O zaman:
  \begin{tagenumerate}{prvEx,prvAll}
    \tagitem{prvEx}{$\Sat{M}{\lexists[x][!A(x)]}$ ancak ve ancak en az
      bir değişken ataması~$s$ için $\Sat{M}{!A(x)}[s]$.}{}
    \tagitem{prvAll}{$\Sat{M}{\lforall[x][!A(x)]}$ ancak ve ancak bütün
      değişken atamaları~$s$ için $\Sat{M}{!A(x)}[s]$.}{}
\end{tagenumerate}
\end{prop}

\begin{proof}
Alıştırma.
\end{proof}

\begin{prob}
\olref[fol][syn][ass]{prop:sat-quant} önermesini ispatlayın.
\end{prob}

\begin{prob}
\DeclareRobustCommand{\VDash}{\mathrel{||}\joinrel\Relbar}
$\Lang L$'nin !!{function}s içermeyen bir dil olduğunu varsayın.
!!{structure}~$\Struct M$, !!a{constant}~$c$ ve $a \in \Domain M$
verildiğinde $\Struct M[a/c]$ ifadesini,~$\Struct M$ ile aynı olup
yalnızca $\Assign{c}{M[a/c]} = a$ olacak biçimde ayrılan !!{structure}
olarak tanımlayın. !!{sentence}s~$!A$ için $\Struct M \VDash !A$
bağıntısını şöyle tanımlayın:
\begin{enumerate}
\tagitem{prvFalse}{%
  \indcase{!A}{\lfalse}{$\Struct{M} \VDash \indfrm$ değil.}}{}

\tagitem{prvTrue}{%
  \indcase{!A}{\ltrue}{$\Struct{M} \VDash \indfrm$.}}{}

\item \indcase{!A}{\Atom{R}{d_1, \dots, d_n}}{$\Struct M \VDash \indfrm$
  ancak ve ancak $\langle \Assign{d_1}{M}, \dots,
  \Assign{d_n}{M} \rangle \in \Assign{R}{M}$.}

\item \indcase{!A}{\eq[d_1][d_2]}{$\Struct M \VDash \indfrm$ ancak ve
  ancak $\Assign{d_1}{M} = \Assign{d_2}{M}$.}

\tagitem{prvNot}{%
  \indcase{!A}{\lnot !B}{$\Struct{M} \VDash \indfrm$ ancak ve ancak
    $\Struct{M} \VDash {!B}$ değil.}}{}

\tagitem{prvAnd}{%
  \indcase{!A}{(!B \land !C)}{$\Struct{M} \VDash
    \indfrm$ ancak ve ancak $\Struct{M} \VDash!B$ ve
    $\Struct{M} \VDash !C$.}}{}

\tagitem{prvOr}{%
  \indcase{!A}{(!B \lor !C)}{$\Struct{M} \VDash \indfrm$ ancak ve ancak
    $\Struct{M} \VDash !B$ ya da $\Struct{M} \VDash !C$
    (veya her ikisi).}}{}

\tagitem{prvIf}{%
  \indcase{!A}{(!B \lif !C)}{$\Struct{M} \VDash \indfrm$ ancak ve ancak
    $\Struct {M} \VDash !B$ değil ya da $\Struct M \VDash !C$
    (veya her ikisi).}}{}

\tagitem{prvIff}{%
  \indcase{!A}{(!B \liff !C)}{$\Struct{M} \VDash {\indfrm}$ ancak ve
    ancak ya hem $\Struct{M} \VDash {!B}$ hem $\Struct{M} \VDash {!C}$,
    ya da ne $\Struct{M} \VDash {!B}$ ne de $\Struct{M} \VDash {!C}$.}}{}

\tagitem{prvAll}{%
  \indcase{!A}{\lforall[x][!B]}{$c$,~$!B$ içinde geçmiyorsa,
    $\Struct{M} \VDash {\indfrm}$ ancak ve ancak her
    $a \in \Domain{M}$ için
    $\Struct{M[a/c]} \VDash \Subst{!B}{c}{x}$.}}{}

\tagitem{prvEx}{%
  \indcase{!A}{\lexists[x][!B]}{$c$,~$!B$ içinde geçmiyorsa,
  $\Struct{M} \VDash {\indfrm}$ ancak ve ancak
  $\Struct{M[a/c]} \VDash \Subst{!B}{c}{x}$ olacak bir
  $a \in \Domain M$ vardır.}}{}
\end{enumerate}
$x_1$, \dots, $x_n$,~$!A$'daki bütün serbest !!{variable}s;
$c_1$, \dots, $c_n$,~$!A$ içinde bulunmayan sabit simgeleri;
$a_1$, \dots, $a_n \in \Domain M$ ve $s(x_i) = a_i$ olsun.

$\Sat{M}{!A}[s]$ ancak ve ancak
$\Struct M[a_1/c_1,\dots,a_n/c_n]
\VDash \Subst{\Subst{!A}{c_1}{x_1}\dots}{c_n}{x_n}$ olduğunu gösterin.

(Bu problem, birinci dereceden mantığa ilişkin semantiğin değişken
atamaları kullanılmadan da verilebileceğini gösterir.)
\end{prob}

\begin{prob}
$f$'nin~$!A(x,y)$ içinde bulunmayan bir fonksiyon simgesi olduğunu
varsayın. $\Sat{M}{\lforall[x][\lexists[y][!A(x,y)]]}$ olacak bir
!!a{structure}~$\Struct{M}$ bulunmasının, ancak ve ancak
$\Sat{M'}{\lforall[x][!A(x,f(x))]}$ olacak bir~$\Struct M'$ bulunmasıyla
eşdeğer olduğunu gösterin.

(Bu problem, Skolem Teoremi diye bilinen sonucun özel bir durumudur;
$\lforall[x][!A(x,f(x))]$,~$\lforall[x][\lexists[y][!A(x,y)]]$ ifadesinin
\emph{Skolem normal biçimi} diye adlandırılır.)
\end{prob}

\end{document}
